\documentclass[11pt]{article}


\usepackage{fullpage}
\usepackage{graphicx}
\usepackage{amsmath}
\usepackage{amssymb}
\usepackage{amsthm}
\usepackage{fancyvrb}

\newcommand{\myname}{Mehshan Mustafa}

\newenvironment{theorem}[2][Theorem]{\begin{trivlist}
\item[\hskip \labelsep {\bfseries #1}\hskip \labelsep {\bfseries #2.}]}{\end{trivlist}}
\newenvironment{lemma}[2][Lemma]{\begin{trivlist}
\item[\hskip \labelsep {\bfseries #1}\hskip \labelsep {\bfseries #2.}]}{\end{trivlist}}
\newenvironment{exercise}[2][Exercise]{\begin{trivlist}
\item[\hskip \labelsep {\bfseries #1}\hskip \labelsep {\bfseries #2.}]}{\end{trivlist}}
\newenvironment{problem}[2][Problem]{\begin{trivlist}
\item[\hskip \labelsep {\bfseries #1}\hskip \labelsep {\bfseries #2.}]}{\end{trivlist}}
\newenvironment{question}[2][Question]{\begin{trivlist}
\item[\hskip \labelsep {\bfseries #1}\hskip \labelsep {\bfseries #2.}]}{\end{trivlist}}
\newenvironment{corollary}[2][Corollary]{\begin{trivlist}
\item[\hskip \labelsep {\bfseries #1}\hskip \labelsep {\bfseries #2.}]}{\end{trivlist}}
\newenvironment{solution}{\begin{proof}[Solution]}{\end{proof}}
\newenvironment{idea}[2][Proof Idea.]{\textit{#1} #2}



\parindent0in
\pagestyle{plain}
\thispagestyle{plain}

\usepackage{csquotes}
\usepackage[shortlabels]{enumitem}

\newcommand{\dated}{\today}
\newcommand{\token}[1]{\langle \text{#1} \rangle}

\begin{document}

\textbf{Introduction to the Theory of
Computation}\hfill\textbf{\myname}\\[0.01in]
\textbf{Chapter 8: Space Complexity}\hfill\textbf{\dated}\\
\smallskip\hrule\bigskip

\begin{problem}{8.21}
For any positive integer $x$, let $x^R$ be the integer whose binary representation is the reverse of the binary representation of $x$ (Assume no leading 0s in the binary representation of $x$.) Define the function $R^+ : \mathbb{N} \longrightarrow \mathbb{N}$ where $R^{+}(x) = x + x^{R}$.
\end{problem}

\begin{problem}[Part]{a}
Let $A_2 = \{\langle x, y \rangle \ | \ R^+(x) = y\}$. Show $A_2 \in$ L.
\end{problem}

\begin{proof}
% Use 3 pointers and a carry bit
% Pointer P1 points at the most significant bit of x
% Pointer P2 points at the least significant bit of x
% Pointer P3 points at the least significant bit of y
% Let n = |x|
% Let c be the carry bit
% Iterate for n - 1 steps:
%	Add P1 + P2 + c, compare result with P3 and update c
%	Increment P1 and P3
%	Decrement P2
To show that $A_2 \in$ L, we give a log space algorithm $M$ that decides $A_2$. \\

$M =$ \textquotedblleft On input $\langle x, y \rangle$, where $x$ and $y$ are positive binary integers without leading 0s:
\begin{enumerate}
\item Let $i$, $j$, $n_x$, $n_y$ be binary integers, and $s$ and $c$ be binary digits on work tape.
\item Let $n_x$ = $|x|$, and $n_y$ = $|y|$.
\item Let $i$ = 1, $j$ = $n_x$, and $c$ = 0.
\item Repeat for $n_x$ steps:
\item \hspace*{0.5cm} If $y$ does not have $i^{th}$ digit, then \textit{reject}.
\item \hspace*{0.5cm} Let $s$ be the sum of $i^{th}$ and $j^{th}$ binary digits of $x$ and $c$. If the addition produces a carry, \\
\hspace*{0.5cm} then $c = 1$. Otherwise, $c = 0$.
\item \hspace*{0.5cm} If $s$ and $i^{th}$ digit of $y$ are different, then \textit{reject}.
\item \hspace*{0.5cm} Increment $i$ by 1, and decrement $j$ by 1.
\item If $c = 1$, then:
\item \hspace*{0.5cm} If $y$ does not have $i^{th}$ digit, then \textit{reject}.
\item \hspace*{0.5cm} If $c$ and $i^{th}$ digit of $y$ are different, then \textit{reject}.
\item If $y$ has $(i+1)^{th}$ digit, then \textit{reject}.
\item \textit{Accept}.\textquotedblright
\end{enumerate}
\end{proof}

\newpage

\begin{problem}[Part]{b}
Let $A_3 = \{\langle x, y \rangle \ | \ R^+(R^+(x)) = y\}$. Show $A_3 \in$ L.
\end{problem}

\begin{proof}
A log space algorithm that decides $A_3$ can be devised like follows:
\begin{enumerate}
\item Compute first and last binary digits of $R^+(x)$ using the algorithm presented in Part a.
\item Add the binary digits and keep track of any carry for later use.
\item Compare the least significant digit of the sum with the least significant digit of $y$. \textit{Reject}, if they are different.
\item Next, iteratively compute the second and second last digits, third and third last digits and so forth.
\item Take sum of the two digits computed in each step and any carry from previous step, and compare the least significant digit of the sum  with the corresponding digit in $y$. \textit{Reject}, if the digits differ.
\item Finally, perform checks related to final carry and extra digits in $y$.
\end{enumerate}
 
\end{proof}

\end{document}